\begin{Solbook}{1}
 \par \noindent  \begin {listEX}[3] \item $(-33) + 45 + (-67)$\\ $= 45 - 33 - 67$\\ $= 12 - 67$\\ $= -55$ \item $28\cdot 69 + 28\cdot 31 + (-100)$\\ $= 28(69 + 31) - 100$\\ $= 28\cdot 100 - 100$\\ $= 2800 - 100$ $= 2700$ \item $75:\left [ (153 - 6\cdot 5^2)\cdot 8 + 2011^0\right ]$\\ $= 75:\left [(153 - 150)\cdot 8 + 1\right ]$\\ $= 75:(3\cdot 8 + 1)$\\ $= 75:25$ $= 3$ \end {listEX}  \hfill \qedEX 
\end{Solbook}
\begin{Solbook}{2}
 \par \noindent  \begin {listEX}[3] \item $6x + 15 = 75$\\ \hspace {2 cm}$\begin {aligned} 6x &= 75 - 15\\ 6x &= 60\\ x &= 10 \end {aligned}$ \item $x - 12 = -27$\\ \hspace {2 cm}$\begin {aligned} x &= -27 + 12\\ x &= -15 \end {aligned}$ \item $23 + (x - 5):4 = 3^8 : 3^5$\\ \hspace {1 cm}$\begin {aligned} 23 + (x - 5):4 &= 3^3\\ 23 + (x - 5):4 &= 27\\ (x - 5):4 &= 4\\ x - 5 &= 16\\ x &= 21 \end {aligned}$ \end {listEX} \par  \hfill \qedEX 
\end{Solbook}
\begin{Solbook}{3}
 \par \noindent  Gọi $x$ là số nhóm nhiều nhất có thể chia $(x\in \mathbb N)$\\ $x =$ ƯCLN$(180;102)$ \par $180=2^2\cdot 3^2\cdot 5\qquad 102=2\cdot 3\cdot 17$ \par ƯCLN$(180;102)=2\cdot 3=6$\\ Vậy có thể chia nhiều nhất 6 nhóm. Mỗi nhóm có $30$ bạn nữ và $17$ bạn nam.  \hfill \qedEX 
\end{Solbook}
\begin{Solbook}{4}
 \par \noindent  \begin {enumerate} \item Diện tích trồng rau là: $30\cdot 12 - 9\cdot 9 = 279$ (m$^2$). \item Số gram hạt giống cần để gieo hết phần đất trồng rau là: $20\cdot 279 = 5\,580$ (g). \end {enumerate}  \hfill \qedEX 
\end{Solbook}
\begin{Solbook}{5}
 \par \noindent  Số tiền thu được khi bán 20 cái máy tính đầu tiên là: $20 \cdot 400\,000 = 8\,000\,000$ (đồng) \par Số tiền thu được khi bán hết số máy tính còn lại là: $20 \cdot 380\,000 = 7\,600\,000$ (đồng) \par Số tiền nhà sách thu được khi bán hết 40 cái máy tính là: $8\,000\,000 + 7\,600\,000 = 15\,600\,000$ (đồng) \par Vì $15\,600\,000$ đồng $> 14\,000\,000$ đồng nên nhà sách lời. \par Số tiền nhà sách lời là: $15\,600\,000 - 14\,000\,000 = 1\,600\,000$ (đồng)  \hfill \qedEX 
\end{Solbook}
\begin{Solbook}{6}
 \par \noindent Để $(3n+7)$ là bội của $(n+1)$ thì $(3n+7)$ phải chia hết cho $(n+1)$.\\ Ta có: $3n+7 = 3(n+1) + 4$ \par Để $3n+7$ chia hết cho $n+1$ thì $4$ phải chia hết cho $n+1$. \par Điều này có nghĩa là $(n+1)$ là ước của $4$. \par $(n+1)\in $Ư$(4)=\{\pm 1, \pm 2, \pm 4\}$. $ \Rightarrow n\in \{-5;-3;-2;0;1;3\}$\\ Nên số nguyên tố $n$ thỏa mãn đề bài là $n = 3$. \hfill \qedEX 
\end{Solbook}
